\chapter[Dynamics and neural manifolds]{Dynamical systems, inverse problems, and neural manifolds}
A state vector $x(t)$ represents a point in a state space.
An ordinary differential equation (ODE) $\dot x=f(x,t)$ defines
how that point moves, provided a solution exists. A vector of
measured rates or LFP values is an observation of the state,
not necessarily the complete physical state itself.

\entry{81}{Picard--Lindelof existence and uniqueness}
{If $f(t,x)$ is locally Lipschitz in $x$ and continuous in $t$
on a suitable neighborhood, a unique local solution exists
through a given initial condition.}
{Determine when deterministic neural state equations define a
well-posed trajectory.}
\entry{82}{Gronwall's inequality}
{A nonnegative function bounded by an integral inequality
has an exponential comparison bound under its hypotheses.}
{Quantify growth of perturbations to trajectories and models.}
\entry{83}{Lyapunov stability theorems}
{A differentiable positive-definite function with suitable
nonpositive or negative orbital derivative implies stability
or asymptotic stability under the corresponding assumptions.}
{Assess recurrent-circuit attractors and locally stable activity.}
\entry{84}{LaSalle's invariance principle}
{In a compact positively invariant region with $\dot V\leq0$,
trajectories approach the largest invariant set where
$\dot V=0$, subject to standard hypotheses.}
{Study convergence when strict energy decrease is unavailable.}
\entry{85}{Hartman--Grobman theorem}
{Near a hyperbolic equilibrium, a smooth autonomous system
is locally topologically conjugate to its linearization.}
{Interpret local fixed-point geometry while respecting nonlinear effects.}
\entry{86}{Kalman observability rank test}
{The linear time-invariant pair $(A,C)$ with $n$ states
is observable exactly when
$\rank[C^\top,(CA)^\top,\ldots,(CA^{n-1})^\top]^\top=n$.}
{Assess recoverability of latent states from neural measurements.}
\entry{87}{Kalman controllability rank test}
{The linear time-invariant pair $(A,B)$ is controllable
exactly when
$\rank[B,AB,\ldots,A^{n-1}B]=n$.}
{Characterize reachable states under modeled input interventions.}
\entry{88}{Kalman filter optimality for linear-Gaussian models}
{With known linear dynamics, Gaussian initial state and
Gaussian white process/measurement noise, Gaussian conditioning
yields recursive conditional means and covariances.}
{Estimate trajectories from noisy population observations.}
\entry{89}{Takens' delay-embedding theorem}
{For a generic smooth observation of a compact $d$-dimensional
manifold and generic diffeomorphism, sufficiently many delay
coordinates (classically $2d+1$) embed its state manifold.}
{Understand the idealized mathematical conditions for reconstructing
dynamics from scalar time series.}
\entry{90}{Whitney embedding theorem}
{Every smooth $d$-dimensional manifold embeds smoothly into
$\R^{2d}$ (with refinements and specific hypotheses).}
{Motivate low-dimensional state representations, not neuronal sampling guarantees.}
\entry{91}{Implicit function theorem}
{If a differentiable constraint $F(x,y)=0$ has an invertible
derivative in $y$ at a point, it defines $y$ locally as a
differentiable function of $x$.}
{Understand local tangent structure of constraint-defined manifolds.}
\entry{92}{Tikhonov regularization convergence results}
{For a bounded linear inverse problem in Hilbert spaces, with
proper regularization conditions and parameter choice relative
to noise, regularized estimates converge to an appropriate
minimum-norm solution as noise vanishes.}
{Describe what stable lead-field inversion can recover and when
source nonuniqueness remains.}

\begin{definition}
A linear state-space model has
\[
x_{t+1}=Ax_t+Bu_t+w_t,\qquad y_t=Cx_t+v_t.
\]
The dimensions, covariance models, initial distribution, and
observation mechanism must be declared before inferential conclusions.
\end{definition}

\begin{caution}
Whitney and Takens do not guarantee that an arbitrary sparse neuron
sample recovers the whole circuit manifold. Observability,
measurement design, noise, and sampling density are independent concerns.
Similarly, a regularized LFP inverse is not a unique physiological
current reconstruction without additional constraints.
\end{caution}
% --- expanded edition (October 2026) ---
\section{Proofs of selected dynamical results}
\begin{proof}[Proof of Result 82 (Gronwall, integral form)]
Let $u$ be continuous and nonnegative with
$u(t)\le a+\int_0^tb\,u(s)\,ds$ for constants $a,b\ge0$. Set
$U(t)=a+\int_0^tb\,u(s)\,ds$. Then $U'=bu\le bU$, so
$(e^{-bt}U)'\le0$ and $U(t)\le ae^{bt}$. Hence $u(t)\le U(t)\le ae^{bt}$.
\end{proof}
\begin{proof}[Proof of Result 81 (Picard--Lindel\"of), global Lipschitz case]
Suppose $|f(t,x)-f(t,y)|\le L|x-y|$. On $C([0,h])$ with the sup norm,
define $(\Phi x)(t)=x_0+\int_0^tf(s,x(s))\,ds$. Then
$\|\Phi x-\Phi y\|_\infty\le Lh\|x-y\|_\infty$, a contraction for $h<1/L$.
By Banach's theorem (Result 24) there is a unique fixed point, which is
the unique solution on $[0,h]$. Repeating from $t=h$ extends it.
Uniqueness on the whole interval also follows from Gronwall: two
solutions satisfy $|x(t)-y(t)|\le\int_0^tL|x-y|\,ds$, so $a=0$ forces
$x=y$.
\end{proof}
Gronwall also quantifies sensitivity: if two initial conditions differ by
$\delta$, the trajectories differ by at most $\delta e^{Lt}$. For a
chaotic network this bound is nearly attained, and predictions of
individual trajectories lose accuracy exponentially, though statistical
properties may still be predictable.

\section{Lyapunov functions for linear systems}
\begin{theorem}
For $\dot x=Ax$, the origin is asymptotically stable iff for some (and
then every) $Q\succ0$ the Lyapunov equation $A^\top P+PA=-Q$ has a
solution $P\succ0$.
\end{theorem}
\begin{proof}[Proof of sufficiency]
Take $V(x)=x^\top Px>0$ for $x\ne0$. Along trajectories,
$\dot V=x^\top(A^\top P+PA)x=-x^\top Qx<0$ for $x\ne0$. With
$\lambda_{\min}(Q)\|x\|^2\le x^\top Qx$ and $V\le\lambda_{\max}(P)\|x\|^2$,
$\dot V\le-cV$ with $c=\lambda_{\min}(Q)/\lambda_{\max}(P)$, so
$V(t)\le V(0)e^{-ct}\to0$, which is Result 83 in quantitative form.
\end{proof}
For stable $A$, the solution is $P=\int_0^\infty e^{A^\top t}Qe^{At}\,dt$.
The same equation with noise appears in stochastic models: for
$dx=Ax\,dt+B\,dW$, the stationary covariance $\Sigma$ solves
$A\Sigma+\Sigma A^\top+BB^\top=0$, which links the dynamics matrix fitted to
neural data to the covariance of its fluctuations.

\section{Observability and controllability proofs}
\begin{proof}[Proof of Result 86]
In discrete time, the noise-free outputs from $x_0$ are
$y_k=CA^kx_0$. Stacking $k=0,\ldots,n-1$ gives $\mathcal Ox_0$ with the
observability matrix $\mathcal O$. If $\rank\mathcal O=n$, $x_0$ is
determined. Conversely, if $\mathcal Ox_0=0$ for some $x_0\ne0$, the
Cayley--Hamilton theorem expresses every $A^k$, $k\ge n$, as a
combination of $I,\ldots,A^{n-1}$, so $CA^kx_0=0$ for all $k$: the
state $x_0$ is invisible forever.
\end{proof}
Controllability (Result 87) is the dual statement for the pair
$(A^\top,B^\top)$. In a recorded population, a latent direction outside
the observable subspace cannot be estimated from that population no
matter how long one records; in a stimulation experiment, a direction
outside the controllable subspace cannot be driven by that stimulus.

\section{Delay embedding in practice}
Takens' theorem (Result 89) guarantees that the map
$x_t\mapsto(y_t,y_{t-\tau},\ldots,y_{t-(m-1)\tau})$ is an embedding for
generic smooth observations when $m\ge2d+1$. Practice adds three
obstacles absent from the theorem: noise, which blurs the reconstructed
manifold at small scales; finite data, which may not cover the
attractor; and the choice of $\tau$, for which common heuristics are the
first zero of the autocorrelation or the first minimum of the mutual
information. Methods that infer causality from embeddings (convergent
cross mapping) assume deterministic dynamics and fail with strong
stochastic forcing or synchrony between the variables.

\section*{Exercises for this chapter}
\begin{exercise}
Show that $\dot x=x^{2/3}$, $x(0)=0$, has two solutions and identify
which hypothesis of Result 81 fails.
\end{exercise}
\begin{exercise}
Solve $A^\top P+PA=-I$ for $A=\diag(-1,-2)$ and verify $P\succ0$.
\end{exercise}
\begin{exercise}
For $A=\begin{psmallmatrix}0&1\\0&0\end{psmallmatrix}$, determine which
of $C=(1,0)$ and $C=(0,1)$ makes the pair observable.
\end{exercise}
\begin{exercise}
Compute the stationary covariance of $dx=-x\,dt+dW$ from the Lyapunov
equation and compare with the Ornstein--Uhlenbeck formula.
\end{exercise}
